\documentclass[11pt]{article}
\usepackage[margin=1in]{geometry}
\usepackage[T1]{fontenc}
\usepackage{lmodern}
\usepackage{amsmath,amssymb}
\usepackage{hyperref}

\title{Sharp Sobolev Extension Constant for Planar Sectors:\\
An Exact Later Answer to the Wedge Asymptotic}
\author{Run fa\_banach\_001, agent\_lane\_14}
\date{11 August 2026}

\begin{document}
\maketitle

\section*{Status and sources}

\textbf{Status:} \texttt{literature\_implied\_answer (full)}.

The original problem is Remark~4.6 on PDF page~15 of Lotoreichik's
arXiv:1306.1524 \cite{Lot}.  The exact answer is Theorem~1 on PDF page~2 of
Zeraoulia's June 2026 preprint \cite{Zer}, deposited at Zenodo under DOI
\href{https://doi.org/10.5281/zenodo.20534573}{10.5281/zenodo.20534573}.
The later paper does not cite \cite{Lot}; the answer relation follows by
matching the definitions and theorem, hence the provenance classification.

\section*{Source problem}

For a planar wedge $\Omega_\varphi$ of opening $0<\varphi\leq\pi$, let
\[
 \mathcal E(\Omega_\varphi)
 =\inf_E\|E\|_{H^1(\Omega_\varphi)\to H^1(\mathbb R^2)},
 \qquad (Ef)|_{\Omega_\varphi}=f,
\]
where
\[
 \|f\|_{H^1(D)}^2=\int_D\bigl(|f|^2+|\nabla f|^2\bigr)\,dx.
\]
The source proves, with $\theta\downarrow0$,
\[
 \sqrt2+\frac{\sqrt2}{32}\theta^2+o(\theta^2)
 \leq \mathcal E(\Omega_{\pi-\theta})
 \leq \sqrt2+\frac{\sqrt2}{4}\theta+o(\theta),
\]
and asks for the exact asymptotics because the displayed bounds even have
different orders.

\section*{Exact answering theorem}

Put $\beta=2\pi-\varphi$.  Theorem~1 of \cite{Zer} computes the exterior and
full-plane sharp constants as
\[
 C_{\mathrm{out}}(\varphi)^2
 =\max\left\{\frac{\varphi}{\beta},\frac{\beta}{\varphi}\right\},
 \qquad
 C_{\mathrm{full}}(\varphi)^2=1+C_{\mathrm{out}}(\varphi)^2
 =\frac{2\pi}{\min\{\varphi,\beta\}}.
\]
Its $C_{\mathrm{full}}$ is exactly the source's $\mathcal E$: both include the
fixed $H^1$ contribution on the original sector as well as the added exterior
part.  Therefore, for $0<\varphi\leq\pi$,
\[
 \boxed{\mathcal E(\Omega_\varphi)=\sqrt{\frac{2\pi}{\varphi}}}.
\]
Substituting $\varphi=\pi-\theta$ gives the complete answer
\[
 \mathcal E(\Omega_{\pi-\theta})
 =\sqrt2\left(1-\frac{\theta}{\pi}\right)^{-1/2}
 =\sqrt2+\frac{\sqrt2}{2\pi}\theta
 +\frac{3\sqrt2}{8\pi^2}\theta^2+O(\theta^3).
\]
In particular, the first correction is linear, with sharp coefficient
$\sqrt2/(2\pi)$.

\section*{Why the sharpness proof matches the source norm}

The upper operator on the complementary sector is the angular
reflection--dilation
\[
 (T_0u)(r,\eta)
 =u\!\left(r,\frac{\varphi(2\pi-\eta)}{\beta}\right),
 \qquad \varphi<\eta<2\pi.
\]
A polar-coordinate change of variables gives
\[
 \|T_0u\|_{H^1(\Omega_\varphi^c)}^2
 =\frac{\beta}{\varphi}
   \int_{\Omega_\varphi}(|u|^2+|u_r|^2)
 +\frac{\varphi}{\beta}
   \int_{\Omega_\varphi}\frac{|u_\vartheta|^2}{r^2},
\]
so the exterior upper constant is the stated maximum.

For sharpness, $t=\log r$ transforms the Dirichlet energy conformally into
the standard energy on an infinite strip.  Slowly varying equal boundary
traces force the ratio $\beta/\varphi$; slowly varying opposite traces force
$\varphi/\beta$.  The Fourier strip calculation is on PDF pages~3--5 of
\cite{Zer}.  Finally, spatial dilation leaves the two-dimensional Dirichlet
energy unchanged while multiplying the $L^2$ term by $\rho^2$; letting
$\rho\downarrow0$ proves that the inhomogeneous $H^1$ constant cannot be
smaller.  Thus no homogeneous/inhomogeneous or exterior/full-plane convention
gap remains.

\begin{thebibliography}{9}
\bibitem{Lot}
V. Lotoreichik,
\emph{Lower bounds on the norms of extension operators for Lipschitz domains},
arXiv:1306.1524 (2013).

\bibitem{Zer}
R. Zeraoulia,
\emph{Sharp $W^{1,2}$ Extension Constants for Planar Sectors},
Zenodo, 4 June 2026, DOI: 10.5281/zenodo.20534573.
\end{thebibliography}

\end{document}
